\chapter{Proposed Method}
\label{chap:method}

As discussed in Chapter~\ref{chap:background}, byte-level fuzzers cannot effectively test database engines because random mutations produce parser-rejected inputs that never reach deep program logic. Grammar-based fuzzers solve the syntactic validity problem but treat all SQL constructs equally, with no knowledge of which structural patterns are most likely to trigger vulnerabilities. This chapter presents DBMS-Nautilus, a grammar-based greybox fuzzing system that combines Nautilus-style coverage-guided generation~\cite{nautilus} with a domain-specific SQL grammar engineered to target vulnerability-triggering code paths in SQLite.

Section~\ref{sec:overview} describes the system architecture. Section~\ref{sec:grammar-design} presents the grammar design methodology, which constitutes the core contribution of this work. Section~\ref{sec:crash-pipeline} covers the crash analysis pipeline.

\section{System Overview}
\label{sec:overview}

DBMS-Nautilus is a closed-loop fuzzing pipeline where coverage feedback from each execution guides subsequent input generation. The system consists of five components, illustrated in Figure~\ref{fig:architecture}: a grammar engine that defines the space of possible SQL inputs, a generation and mutation module that constructs derivation trees, an execution layer (fork server and harness) that runs each input against an instrumented SQLite build, a coverage feedback mechanism that identifies inputs exercising new code paths, and a triage pipeline that analyzes crashes after each campaign.

\begin{figure}[htbp]
\centering
\begin{tikzpicture}[
    node distance=1.4cm,
    comp/.style={rectangle, draw, rounded corners, minimum width=3.2cm, minimum height=1.0cm, text centered, align=center, font=\small},
    arrow/.style={->, >=stealth, thick},
    lbl/.style={font=\scriptsize, fill=white, inner sep=1pt}
]
    % Components in a flow
    \node[comp] (grammar) {\textbf{1} Grammar Engine\\(Python DSL + PyO3)};
    \node[comp, right=2.0cm of grammar] (gen) {\textbf{2} Generation\\+ Mutation};
    \node[comp, below=of gen] (exec) {\textbf{3} Fork Server\\+ SQLite Harness};
    \node[comp, left=2.0cm of exec] (feedback) {\textbf{4} Coverage\\Feedback + Queue};
    \node[comp, below=1.8cm of feedback] (triage) {\textbf{5} Triage Pipeline};

    % Main loop
    \draw[arrow] (grammar) -- node[lbl, above] {rules + weights} (gen);
    \draw[arrow] (gen) -- node[lbl, right] {SQL string (file)} (exec);
    \draw[arrow] (exec) -- node[lbl, below] {bitmap + exit status} (feedback);
    \draw[arrow] (feedback) -- node[lbl, left] {interesting inputs} (gen);

    % Triage branch
    \draw[arrow] (exec.south) -- ++(0,-0.5) -| node[lbl, below left] {crashes} (triage.east);
\end{tikzpicture}
\caption{Architecture of DBMS-Nautilus. Components 1--4 form a closed feedback loop: the grammar defines the input space, generation constructs SQL from derivation trees, the fork server executes each input against an instrumented SQLite harness, and coverage feedback retains inputs that discover new edges. Component 5 analyzes crashes after each campaign.}
\label{fig:architecture}
\end{figure}

\paragraph{Grammar Engine.}
The grammar engine loads a SQL grammar specification written in Python via the PyO3 foreign function interface. Each grammar file defines production rules using the \texttt{ctx.rule()} API, where non-terminals are enclosed in curly braces (e.g., \texttt{ctx.rule("Stmt", "SELECT \{Expr\} FROM \{Table\}")}). The engine builds a rule table from the parsed grammar and supports weighted sampling over production rules, allowing certain rules to be selected more frequently than others (Section~\ref{sec:grammar-design} describes how weights are assigned).

\paragraph{Generation and Mutation.}
On each iteration, the fuzzer either generates a fresh derivation tree from the grammar or selects an existing tree from the queue and applies a mutation. Three tree-level mutations are used: random mutation (replace a subtree with a freshly generated one for the same non-terminal), random recursion (expand or collapse recursive productions), and splice mutation (swap compatible subtrees between two queue trees). The resulting tree is unparsed into a concrete SQL string by reading its leaf nodes left to right.

\paragraph{Fork Server and Harness.}
The SQL string is written to a temporary file and passed as a command-line argument to the SQLite harness. The fork server, following the AFL protocol~\cite{afl}, loads the harness once into a parent process and creates child processes via \texttt{fork()} for each test case. Each child inherits the initialized program state, reads the SQL file, executes it via \texttt{sqlite3\_exec()}, and exits. The parent collects the child's exit status and reads the shared memory coverage bitmap. This fork-based execution eliminates the overhead of process creation on every test case, achieving thousands of executions per second. SQLite is compiled from source using \texttt{afl-clang-fast} with AddressSanitizer and UndefinedBehaviorSanitizer instrumentation, enabling the detection of memory safety violations and undefined behavior during execution.

\paragraph{Coverage Feedback and Queue.}
After each execution, the fuzzer compares the resulting coverage bitmap against the global coverage map accumulated from all prior executions. The bitmap is a shared memory region where each byte represents a hashed edge (basic block transition) in SQLite's control flow graph. If the input exercised at least one previously unseen edge, it is classified as interesting and added to the queue. This queue serves as the seed pool for subsequent mutations, creating a directed search that progressively steers generation toward unexplored program regions. This feedback loop is what makes DBMS-Nautilus a \textit{greybox} fuzzer rather than a purely generative one.

\paragraph{Triage Pipeline.}
After a campaign completes, the triage pipeline processes all collected crashes. It deduplicates crashes by stack hash, matches them against known CVE signatures, and scores reproduction fidelity. Section~\ref{sec:crash-pipeline} describes this pipeline in detail.

A key design decision in DBMS-Nautilus is that the harness starts with an empty in-memory database. Unlike SQLsmith~\cite{sqlsmith}, which connects to an existing database and introspects its schema to generate valid queries, or Squirrel~\cite{squirrel}, which maintains a schema model during fuzzing, DBMS-Nautilus requires no external database state. Instead, the grammar itself generates complete, self-contained SQL that includes schema creation (via the \texttt{Schema-Setup} non-terminal) before issuing queries. Each test case creates its own tables, inserts its own data, and runs its own queries --- making every execution independently reproducible without any setup beyond the empty database. This design simplifies the harness, eliminates dependencies on pre-loaded state, and ensures that any crashing input can be replayed in isolation.

Listing~\ref{lst:harness} shows the harness implementation. The \texttt{setup\_db()} constructor opens the in-memory database before the \texttt{\_\_AFL\_INIT()} fork point, so every child process inherits the initialized (but empty) database handle. Each child reads one SQL file, executes it, and exits.

\begin{lstlisting}[language=C, caption={SQLite harness for AFL fork-server mode. The database is opened before the fork point; each child executes one SQL input from a file.}, label=lst:harness, basicstyle=\ttfamily\small, numbers=left]
#include "sqlite3.h"
static sqlite3 *db = NULL;

__attribute__((constructor))
static void setup_db(void) {
    sqlite3_open(":memory:", &db);
}

int main(int argc, char **argv) {
    __AFL_INIT();  /* Fork point */
    char *sql = read_input(argv[1]);
    if (sql) {
        char *errmsg = NULL;
        sqlite3_exec(db, sql, NULL, NULL, &errmsg);
        if (errmsg) sqlite3_free(errmsg);
        free(sql);
    }
    sqlite3_close(db);
    return 0;
}
\end{lstlisting}

Table~\ref{tab:components} summarizes the five components, their implementation languages, and their roles in the pipeline.

\begin{table}[htbp]
\caption{DBMS-Nautilus system components.}
\label{tab:components}
\centering
\begin{tabular}{llll}
\toprule
\textbf{Component} & \textbf{Language} & \textbf{Module} & \textbf{Role} \\
\midrule
Grammar Engine & Rust & \texttt{grammartec/} & Load grammar, weighted rule sampling \\
Generation + Mutation & Rust & \texttt{fuzzer/} & Build and mutate derivation trees \\
Fork Server + Harness & Rust + C & \texttt{forksrv/}, \texttt{harness/} & Execute SQL against instrumented SQLite \\
Coverage Feedback & Rust & \texttt{fuzzer/} & Track edges, maintain queue \\
Triage Pipeline & Python & \texttt{triage/} & Crash dedup, CVE matching, reporting \\
\bottomrule
\end{tabular}
\end{table}

%% ====================================================================
\section{Structural Primitives Philosophy}
\label{sec:primitives-philosophy}
%% ====================================================================

The grammar encodes SQL structural patterns that trigger vulnerability classes rather than specific proof-of-concept inputs. This ``zero PoC contamination'' principle ensures that the fuzzer discovers bugs through compositional exploration rather than replaying known attacks. For example, the CVE-2020-13434 integer overflow~\cite{cve202013434} requires both a \texttt{printf} function call and an INT32 boundary value. Rather than encoding the exact PoC \texttt{SELECT printf('\%.*g', 2147483647, 0.01)}, the grammar provides independent non-terminals for format specifiers (\texttt{Format-Spec}), boundary integers (\texttt{Boundary-Int}), and function call syntax (\texttt{Boundary-Func-Call}). The fuzzer must compose these building blocks into a triggering combination through random generation and mutation, demonstrating genuine discovery capability.

This philosophy yields three concrete design constraints:

\begin{enumerate}[label=(\alph*)]
    \item No grammar rule may contain a complete CVE proof-of-concept as a fixed string.
    \item Every structural pattern needed to reach a target CVE must exist as a composable non-terminal.
    \item The grammar must preserve sufficient compositional freedom that the fuzzer can discover novel combinations beyond the known CVEs.
\end{enumerate}

Constraint~(a) prevents the fuzzer from trivially replaying known attacks, ensuring that any CVE rediscovery constitutes a genuine compositional achievement. Constraint~(b) guarantees that the grammar's expressible input space contains the structural ingredients required for each target CVE --- without these building blocks, no amount of mutation could reach the vulnerable code path. Constraint~(c) ensures that the grammar is not so narrowly tailored to known CVEs that it cannot discover new bugs: the non-terminals must combine freely in ways the grammar designer did not explicitly anticipate.

The tension between constraints (b) and (c) is the central design challenge. Too few non-terminals and the fuzzer lacks the vocabulary to reach target vulnerabilities. Too many non-terminals and the input space becomes so large that the probability of assembling the right combination within a finite campaign drops below a useful threshold. The grammar evolution process described in Section~\ref{sec:grammar-evolution} addresses this tension through iterative refinement guided by experimental feedback.

%% ====================================================================
\section{Layer Decomposition}
\label{sec:layer-decomposition}
%% ====================================================================

The grammar employs a two-layer architecture that separates atomic SQL constructs from composed vulnerability-triggering patterns.

\subsection{Layer 1: SQL Atoms}

Layer~1 defines the SQL primitives: over 30 expression alternatives (arithmetic, comparison, logical, string, subquery), SELECT/FROM/WHERE/JOIN clauses, window functions with frame specifications, common table expressions, DML statements (INSERT, UPDATE, DELETE), DDL statements (CREATE TABLE, VIEW, INDEX, TRIGGER), function calls (aggregate, scalar, window-specific), and terminal literals (integers, floats, strings, blobs, boundary values). These constitute the atomic building blocks of SQL.

\subsection{Layer 2: Composed Shapes}

Layer~2 composes Layer~1 primitives into four high-level shapes. The \texttt{Schema-Setup} non-terminal provides six schema construction alternatives. The \texttt{Stress-Query} non-terminal provides eight query patterns designed to exercise complex query planning paths. The \texttt{Validation-Op} non-terminal provides four verification operations that check internal database consistency. The \texttt{Boundary-Func-Call} non-terminal provides function calls with boundary values targeting integer overflow and buffer overread conditions.

Listing~\ref{lst:schema-setup} shows the Layer~2 Schema-Setup alternatives. Each alternative creates a different database schema topology that enables distinct classes of vulnerability-triggering queries. The weights assigned to each rule reflect the empirical importance of each schema pattern for triggering bugs: the generated-column schema (S2) receives the highest weight of 3.0 because generated columns are involved in two of the six target CVEs.

\begin{lstlisting}[language=Python, caption={Schema-Setup alternatives (Layer 2). Each alternative constructs a different schema topology. Weights reflect empirical importance for vulnerability triggering.}, label=lst:schema-setup, basicstyle=\ttfamily\small, numbers=left]
# S1: Single table, basic columns
ctx.rule("Schema-Setup",
    "CREATE TABLE IF NOT EXISTS p({Col-Def-List})",
    weight=1.5)

# S2: Table with generated column + constraints
ctx.rule("Schema-Setup",
    "CREATE TABLE IF NOT EXISTS p({Col-Def-List-GenCol})",
    weight=3.0)

# S3: Two tables (enables JOINs between p and q)
ctx.rule("Schema-Setup",
    "CREATE TABLE IF NOT EXISTS p({Col-Def-List});\n"
    "CREATE TABLE IF NOT EXISTS q({Col-Def-List})",
    weight=2.5)

# S4: Table + VIEW
ctx.rule("Schema-Setup",
    "CREATE TABLE IF NOT EXISTS p({Col-Def-List});\n"
    "CREATE VIEW IF NOT EXISTS v1 AS {Select-Stmt}",
    weight=2.0)

# S5: Virtual table (FTS5/FTS3) -- reduced in v3.1
ctx.rule("Schema-Setup",
    "CREATE VIRTUAL TABLE IF NOT EXISTS fts_t1 "
    "USING {Fts-Engine}({Col-Name-List})",
    weight=0.5)

# S6: Table + INDEX
ctx.rule("Schema-Setup",
    "CREATE TABLE IF NOT EXISTS p({Col-Def-List});\n"
    "CREATE INDEX IF NOT EXISTS idx1 ON p({Col-Name})",
    weight=1.0)
\end{lstlisting}

Listing~\ref{lst:stress-query} shows representative Stress-Query alternatives. The NATURAL JOIN pattern (Q3) is particularly relevant because CVE-2020-13435~\cite{cve202013435} requires a NATURAL JOIN in combination with window functions and a coalesce call. The compound query pattern (Q5) exercises the INTERSECT and EXCEPT operators needed for CVE-2020-15358~\cite{cve202015358} and CVE-2020-13871~\cite{cve202013871}.

\begin{lstlisting}[language=Python, caption={Representative Stress-Query alternatives (Layer 2). Each pattern exercises a distinct query planning path in the SQLite optimizer.}, label=lst:stress-query, basicstyle=\ttfamily\small, numbers=left]
# Q2: EXISTS subquery
ctx.rule("Stress-Query",
    "SELECT {Result-Col-List} FROM {Table-Name} "
    "WHERE EXISTS ({Select-Stmt})",
    weight=3.0)

# Q3: NATURAL JOIN
ctx.rule("Stress-Query",
    "SELECT {Result-Col-List} FROM {Table-Name} "
    "NATURAL JOIN {Table-Name} WHERE {Expr}",
    weight=3.0)

# Q5: Compound query (INTERSECT/EXCEPT)
ctx.rule("Stress-Query",
    "{Select-Stmt} {Compound-Op} {Select-Stmt}",
    weight=2.5)

# Q6: Self-JOIN + expression
ctx.rule("Stress-Query",
    "SELECT {Result-Col-List} FROM {Table-Name} "
    "JOIN {Table-Name} {Col-Alias} ON {Expr}",
    weight=2.0)
\end{lstlisting}

Listing~\ref{lst:boundary-func} shows the Boundary-Func-Call non-terminal. The \texttt{printf} with format specifier \texttt{'\%.*g'} and boundary integer 2147483647 directly targets the CVE-2020-13434 integer overflow in SQLite's printf implementation. The \texttt{substr} with boundary integers can trigger heap buffer overreads, while \texttt{hex(zeroblob(N))} with large $N$ values exercises memory allocation edge cases.

\begin{lstlisting}[language=Python, caption={Boundary-Func-Call alternatives targeting integer overflow and buffer boundary conditions.}, label=lst:boundary-func, basicstyle=\ttfamily\small, numbers=left]
ctx.rule("Boundary-Func-Call",
    "printf({Format-Spec}, {Boundary-Int}, {Boundary-Float})",
    weight=3.0)
ctx.rule("Boundary-Func-Call",
    "substr({Str-Literal}, {Boundary-Int}, {Boundary-Int})",
    weight=1.5)
ctx.rule("Boundary-Func-Call",
    "hex(zeroblob({Boundary-Int}))", weight=2.0)

ctx.rule("Format-Spec", "'%.*g'", weight=3.0)
ctx.rule("Format-Spec", "'%.*f'", weight=2.0)

ctx.rule("Boundary-Int", "2147483647", weight=3.0)   # INT32_MAX
ctx.rule("Boundary-Int", "2147483648", weight=3.0)   # INT32_MAX+1
ctx.rule("Boundary-Int", "-2147483648", weight=3.0)  # INT32_MIN
ctx.rule("Boundary-Int", "9223372036854775807", weight=3.0) # INT64_MAX
\end{lstlisting}

%% ====================================================================
\section{Grammar Evolution}
\label{sec:grammar-evolution}
%% ====================================================================

The grammar underwent iterative refinement across four versions, guided by experimental feedback from fuzzing campaigns. Each version addresses a specific deficiency identified through crash analysis, coverage measurement, or CVE reachability assessment. Table~\ref{tab:grammar-evolution} summarizes the version history.

\begin{table}[htbp]
\caption{Grammar version history. Each version addresses a specific deficiency identified through experimental feedback.}
\label{tab:grammar-evolution}
\centering
\begin{tabular}{llrl}
\toprule
\textbf{Version} & \textbf{Date} & \textbf{Rules} & \textbf{Key Change} \\
\midrule
v3.0 & 2026-04-27 & 449 & Initial structural primitives design \\
v3.1 & 2026-04-28 & 449 & S5 (FTS) weight 2.0$\to$0.5; diversify crash portfolio \\
v3.2 & 2026-05-05 & 475 & +window functions, +self-ref generated columns, \\
     &            &     & +count() zero-arg, +ORDER BY 3-term \\
v3.3 & 2026-05-10 & 520 & +JSON/JSONB expansion (Json-Key, Json-Path, Json-Literal) \\
\bottomrule
\end{tabular}
\end{table}

\subsection{v3.0: Initial Design (449 Rules)}

Version~3.0 established the structural primitives architecture described in Sections~\ref{sec:primitives-philosophy} and~\ref{sec:layer-decomposition}. The grammar comprised 449 production rules organized into the two-layer hierarchy, with Layer~1 providing SQL atoms and Layer~2 composing them into schema-setup, stress-query, validation, and boundary function call shapes. Initial weights were assigned based on domain analysis of which structural patterns are associated with each target CVE. This version served as the baseline for subsequent refinement.

\subsection{v3.0 to v3.1: Weight Rebalancing (449 Rules)}

The first refinement addressed a crash diversity problem. In campaigns run with v3.0, 92\% of crashes involved FTS (Full-Text Search) virtual tables. The root cause was straightforward: the S5 schema variant (FTS virtual table creation) had weight 2.0, while the average weight across other schema variants was approximately 1.5. Because the FTS module in SQLite contains many shallow bugs --- particularly in its tokenizer and auxiliary function implementations --- even moderate overweighting of FTS schemas caused FTS-related crashes to dominate the crash portfolio.

The fix was to reduce the S5 weight from 2.0 to 0.5, making FTS schemas the least likely schema variant. This single weight change diversified the crash portfolio significantly: subsequent campaigns produced crashes from generated columns, views, joins, and index operations that had previously been crowded out by FTS bugs. Version~3.1 retained the same 449 rules; only the weight distribution changed.

This episode illustrates a general principle of grammar engineering for fuzzing: a grammar that is structurally capable of reaching a vulnerability may still fail to reach it in practice if the weight distribution concentrates sampling effort on a different, shallower region of the bug space. Weight tuning is a necessary complement to structural coverage.

\subsection{v3.1 to v3.2: Structural Expansion (475 Rules)}

The second refinement addressed structural gaps that prevented the grammar from reaching four of the six target CVEs. A systematic gap analysis compared the structural requirements of each CVE (as documented in Table~\ref{tab:cve-signatures} in Section~\ref{sec:triage-cve}) against the non-terminals available in v3.1. The analysis revealed that:

\begin{itemize}
    \item CVE-2020-9327 required self-referential generated columns (a generated column whose expression references another generated column in the same table), which v3.1 could not express.
    \item CVE-2020-13435 required window functions with \texttt{OVER} clauses, which v3.1 lacked entirely.
    \item CVE-2020-13871 required window functions combined with multi-column \texttt{ORDER BY} clauses.
    \item CVE-2020-15358 required compound queries with \texttt{INTERSECT}/\texttt{EXCEPT} operators, which were present but lacked the specific \texttt{ORDER BY} clause structure needed.
\end{itemize}

Version~3.2 added 26 new rules to close these gaps: 15 window function rules covering \texttt{lead}, \texttt{lag}, \texttt{row\_number}, \texttt{rank}, \texttt{dense\_rank}, \texttt{ntile}, \texttt{first\_value}, \texttt{last\_value}, \texttt{nth\_value}, \texttt{percent\_rank}, and \texttt{cume\_dist} with \texttt{OVER} clause variants (including \texttt{PARTITION BY} and frame specifications); 3 self-referential generated column rules where the generated expression references a sibling column; 1 zero-argument \texttt{count()} form (distinct from \texttt{count(expr)}); 1 three-term \texttt{ORDER BY} clause; and 6 supporting non-terminals for window frame boundaries and partition specifications.

After this expansion, all six target CVEs became structurally reachable --- the grammar contained every building block needed to compose a triggering input for each CVE. CVE reachability improved from 2/6 to 6/6.

\subsection{v3.2 to v3.3: JSON/JSONB Expansion (520 Rules)}

The third refinement targeted the \texttt{json.c} module in SQLite, motivated by the goal of extending the fuzzer's reach beyond the original six CVEs into a current SQLite release (version 3.53.0). The JSON1 extension is a high-value target because it performs complex string parsing, type coercion, and memory management operations that are historically prone to vulnerabilities.

Version~3.3 added 51 new rules organized into four non-terminal groups:

\begin{itemize}
    \item \texttt{Json-Key}: 4 rules for JSON key access patterns (dot notation, bracket notation, quoted keys, nested paths).
    \item \texttt{Json-Path}: 8 rules for JSON path expressions (\texttt{\$} root, array indexing, recursive descent, wildcard).
    \item \texttt{Json-Literal}: 6 rules for JSON value construction (objects, arrays, strings, numbers, booleans, null, nested structures).
    \item JSON function calls: rules for \texttt{json()}, \texttt{json\_extract()}, \texttt{json\_set()}, \texttt{json\_insert()}, \texttt{json\_remove()}, \texttt{json\_type()}, \texttt{json\_array()}, and \texttt{json\_object()}, each composed with the above non-terminals.
\end{itemize}

\subsection{Key Insight: Two Axes of Grammar Evolution}

Grammar evolution operates on two orthogonal axes. \emph{Weight tuning} controls the probability distribution over existing rules, determining which regions of the expressible input space receive more sampling effort. \emph{Structural expansion} extends the boundaries of the expressible input space itself by adding new non-terminals and production rules.

Both axes are necessary. Weight tuning without structural coverage leaves CVEs unreachable regardless of campaign duration --- if the grammar cannot express window functions, no weight adjustment will produce an input containing \texttt{OVER} clauses. Conversely, structural expansion without weight tuning can lead to dominance of one grammar region, as demonstrated by the FTS saturation in v3.0. Effective grammar engineering requires alternating between these two modes of refinement: expanding structure when gap analysis reveals missing building blocks, and tuning weights when crash analysis reveals imbalanced sampling.

%% ====================================================================
\section{Harness Construction}
\label{sec:harness}
%% ====================================================================

\subsection{AFL Fork Server Protocol}

The harness uses the \texttt{\_\_AFL\_INIT()} macro from AFL~\cite{afl} to implement fork-server mode. This is distinct from AFL's persistent mode (\texttt{\_\_AFL\_LOOP}), which Nautilus does not support. In fork-server mode, the parent process loads the SQLite amalgamation into memory, initializes the in-memory database, and then waits for instructions from the fuzzer. For each test case, the parent forks a child process. The child inherits the pre-initialized state, reads the SQL input file, executes it via \texttt{sqlite3\_exec}, and exits. The parent collects the child's exit status and coverage bitmap, then reports results back to the fuzzer coordinator.

The database starts empty because the grammar generates self-contained multi-statement SQL that creates its own schema via the \texttt{Schema-Setup} non-terminal before issuing queries. This design eliminates the need for a pre-loaded schema in the harness itself, simplifying the harness code and ensuring that every test case is independently reproducible.

Listing~\ref{lst:harness} shows a simplified version of the harness implementation. The \texttt{setup\_db} constructor opens the in-memory database before \texttt{\_\_AFL\_INIT()} establishes the fork point. Each forked child reads one SQL file and executes it, then exits.

\begin{lstlisting}[language=C, caption={Simplified SQLite harness for AFL fork-server mode. The database is opened before the fork point; each child executes one SQL input.}, label=lst:harness, basicstyle=\ttfamily\small, numbers=left]
#include "sqlite3.h"

static sqlite3 *db = NULL;

__attribute__((constructor))
static void setup_db(void) {
    sqlite3_open(":memory:", &db);
}

int main(int argc, char **argv) {
    __AFL_INIT();  /* Fork point */

    size_t sql_len = 0;
    char *sql = read_input(argv[1], &sql_len);
    if (sql) {
        char *errmsg = NULL;
        sqlite3_exec(db, sql, NULL, NULL, &errmsg);
        if (errmsg) sqlite3_free(errmsg);
        free(sql);
    }
    sqlite3_close(db);
    return 0;
}
\end{lstlisting}

\subsection{Compilation and Instrumentation}

The harness is compiled using \texttt{afl-clang-fast} with AddressSanitizer (ASan) and UndefinedBehaviorSanitizer (UBSan) enabled. The SQLite amalgamation is compiled as part of the same translation unit, ensuring full instrumentation of all SQLite internals. Additional compile-time flags enable optional SQLite features: \texttt{SQLITE\_ENABLE\_FTS3}, \texttt{SQLITE\_ENABLE\_FTS5}, \texttt{SQLITE\_ENABLE\_RTREE}, \texttt{SQLITE\_ENABLE\_JSON1}, and \texttt{SQLITE\_DEBUG}. The \texttt{SQLITE\_DEBUG} flag enables internal assertion checks that fire \texttt{SIGTRAP} on invariant violations, providing an additional oracle beyond sanitizer-detected errors.

\subsection{Three Harness Types}

Three harness variants are compiled for different stages of the analysis pipeline:

\begin{itemize}
    \item \textbf{\texttt{afl/}} --- the primary fuzzing harness, compiled with \texttt{afl-clang-fast} for AFL coverage instrumentation. Used during fuzzing campaigns.
    \item \textbf{\texttt{test/}} --- the triage harness, compiled with ASan and UBSan but without AFL instrumentation. Used to replay and classify crashes during the triage pipeline.
    \item \textbf{\texttt{nosanit/}} --- the production-mode harness, compiled without any sanitizer instrumentation. Used to assess whether a crash is exploitable under realistic (non-sanitized) conditions.
\end{itemize}

%% ====================================================================
\section{Oracle Classification}
\label{sec:oracle}
%% ====================================================================

The harness produces different exit signals depending on the type of error detected. Table~\ref{tab:oracle} defines the classification scheme. ASan violations (buffer overflows, use-after-free, double-free) cause the process to exit with code 223, configured via the \texttt{ASAN\_OPTIONS=exitcode=223} environment variable. UBSan violations (integer overflow, null pointer dereference, shift out of bounds) cause exit code 1, configured via \texttt{UBSAN\_OPTIONS=halt\_on\_error=1,exitcode=1}. SQLite debug assertions fire \texttt{SIGTRAP} (signal~5), which the fork server detects as exit code 133. Normal execution and SQL semantic errors both produce exit code 0 and are not classified as bugs.

\begin{table}[htbp]
\caption{Oracle classification of harness exit signals. Each signal type corresponds to a distinct vulnerability class.}
\label{tab:oracle}
\centering
\begin{tabular}{llll}
\toprule
\textbf{Signal} & \textbf{Exit Code} & \textbf{Classification} & \textbf{Action} \\
\midrule
ASan violation & 223 & Memory error & Save to \texttt{crashes/} \\
UBSan violation & 1 & Undefined behavior & Save to \texttt{crashes/} \\
SIGTRAP (signal 5) & 133 & Debug assertion & Save to \texttt{signaled/} \\
Normal exit & 0 & No bug & Discard \\
SQL error & 0 & Semantic error & Ignore \\
Timeout & --- & Hang & Save to \texttt{timeout/} \\
\bottomrule
\end{tabular}
\end{table}

%% ====================================================================
\section{Triage Pipeline}
\label{sec:triage}
%% ====================================================================

The triage pipeline transforms raw crash data from a fuzzing campaign into actionable vulnerability reports. A typical hour-long campaign may produce hundreds of crashing inputs, many of which trigger the same underlying bug. The pipeline performs four stages: stack-hash deduplication, CVE signature matching, fidelity scoring, and crash minimization.

\subsection{Stack-Hash Deduplication}

Stack-hash deduplication identifies unique root causes among the collected crashes. For each crashing input, the pipeline replays the input under GDB with the \texttt{test/} harness binary, captures the top 5 stack frames at the crash point, filters out frames belonging to system libraries (\texttt{libc}, \texttt{libasan}, \texttt{libubsan}, \texttt{libpthread}) and the harness wrapper, and retains only SQLite-internal frames. Each retained frame is encoded as a \texttt{function:file:line} triple. The pipeline then computes a SHA-256 hash of the concatenated filtered frames. Crashes that produce identical hashes are considered manifestations of the same root cause and are grouped into a single cluster.

The filtering step is essential because sanitizer runtime frames and C library frames appear in every crash regardless of root cause. By removing these frames, the algorithm focuses on the SQLite-specific call stack that distinguishes one bug from another. The use of the top 5 frames rather than the full stack provides a balance between precision (distinguishing genuinely different bugs) and recall (grouping diverse triggers of the same underlying fault).

Algorithm~\ref{alg:dedup} formalizes the deduplication procedure. The output is a set of clusters, each containing a representative crash, its top frame, and the count of duplicate crashes sharing that root cause.

\begin{algorithm}[htbp]
\caption{Stack-Hash Crash Deduplication}
\label{alg:dedup}
\begin{algorithmic}[1]
\Require Set of crash files $C = \{c_1, \ldots, c_n\}$, harness binary $H$
\Ensure Set of unique root-cause clusters
\For{each crash file $c_i \in C$}
    \State $\text{frames} \gets \Call{RunGDB}{H, c_i}$ \Comment{Top 5 frames}
    \State $\text{filtered} \gets \Call{FilterFrames}{\text{frames}}$ \Comment{Remove libc, ASan, harness}
    \State $\text{hash}_i \gets \text{SHA-256}(\text{Join}(\text{filtered}))$
\EndFor
\State $\text{clusters} \gets \Call{GroupBy}{\{(\text{hash}_i, c_i)\}_{i=1}^n}{\text{hash}}$
\State \Return clusters sorted by decreasing size
\end{algorithmic}
\end{algorithm}

\subsection{CVE Signature Matching}
\label{sec:triage-cve}

CVE signature matching determines whether a crashing input corresponds to a known vulnerability. Each target CVE has a signature defined as a conjunction of structural regex patterns. A crash matches a CVE signature if and only if every required pattern finds at least one occurrence in the SQL input text. The patterns encode structural requirements rather than byte-level identity, using case-insensitive matching with flexible whitespace.

For example, CVE-2020-13434~\cite{cve202013434} requires two structural patterns: a \texttt{printf} function call and an INT32 boundary value (2147483647 or 2147483648). A crash input matches this signature if it contains both \texttt{printf(} and a number matching the regex \texttt{-?214748364[7-8]}. CVE-2020-9327~\cite{cve20209327} requires four structural patterns: a GENERATED column definition, a \texttt{coalesce} function call, a JOIN or comma-join, and a CREATE VIEW statement. Only inputs containing all four patterns simultaneously are classified as CVE-2020-9327 matches.

Table~\ref{tab:cve-signatures} summarizes the required structural patterns for each target CVE. The pattern count varies from 2 (CVE-2019-19646, CVE-2020-13434) to 5 (CVE-2020-13435), reflecting the complexity of the conditions needed to trigger each vulnerability.

\begin{table}[htbp]
\caption{CVE signature patterns. Each CVE requires all listed structural patterns to be present simultaneously in the SQL input for a match.}
\label{tab:cve-signatures}
\centering
\begin{tabular}{lrl}
\toprule
\textbf{CVE} & \textbf{Patterns} & \textbf{Required Structural Elements} \\
\midrule
CVE-2020-15358 & 3 & CREATE VIEW with ORDER BY, INTERSECT in WHERE, \\
               &   & implicit JOIN (comma-separated FROM) \\
CVE-2020-13871 & 3 & EXCEPT, multi-column ORDER BY, scalar subquery \\
CVE-2020-13435 & 5 & NATURAL JOIN, coalesce(), window OVER(), \\
               &   & UNIQUE column, IN subquery \\
CVE-2020-13434 & 2 & printf() call, INT32 boundary value \\
CVE-2020-9327  & 4 & GENERATED column, coalesce(), JOIN, CREATE VIEW \\
CVE-2019-19646 & 2 & GENERATED column, PRAGMA integrity\_check \\
\bottomrule
\end{tabular}
\end{table}

\subsection{Fidelity Scoring}

Not all crashes are deterministically reproducible. Heap layout non-determinism, timing-dependent behavior, and ASLR can cause a crash to reproduce intermittently. The fidelity scoring stage replays each unique crash (after deduplication) a fixed number of times $N$ and records what fraction of replays reproduce the crash. Crashes that reproduce at least 80\% of the time are classified as high-fidelity and are prioritized for further analysis and CVE confirmation. Low-fidelity crashes (below 80\%) may indicate environment-dependent issues or false positives from sanitizer instrumentation.

\subsection{Crash Minimization}

The crash minimization stage iteratively removes SQL statements from a crashing input while preserving the crash behavior. Starting from the full multi-statement input, the algorithm removes one statement at a time and replays the reduced input. If the crash still reproduces, the removal is kept; otherwise, it is reverted. This process continues until no further statements can be removed without losing the crash. The result is a minimal reproducer that contains only the SQL statements necessary to trigger the vulnerability, which is useful for root-cause analysis and CVE confirmation.
